\chapter{Fisher information and the Cram\'er--Rao bound}
\label{ch:fisher-cramer-rao}

The paper's information theorems are statements about Fisher information: how much
one noised sequence tells an estimator about a coordinate of the clean law, and how
that amount decays with the noise level. Its learning section then compares an
estimator's dispersion with the Cram\'er--Rao reference built from that information,
and its corollary on sample complexity invokes the local asymptotic minimax theorem
of H\'ajek and Le Cam. This chapter defines the statistical score and keeps it apart
from the diffusion score, derives Fisher information and its properties, proves the
Cram\'er--Rao bound, works out the efficient information under nuisance parameters and
why a nuisance at another Hermite degree costs nothing to leading order, derives the
Gaussian information formula the paper uses in closed form, states local asymptotic
normality and the H\'ajek--Le Cam theorem in plain words and in the form the paper's
corollary uses, explains what ``maximum likelihood attains it'' does and does not say,
and derives Fisher's identity for latent-variable models, which is how the paper
computes the gradient of its marginal likelihood from one forward--backward pass.

\section{Two scores}
\label{sec:two-scores}

The word \term{score} has two meanings in this project, and confusing them makes the
paper unreadable.

\begin{definition}[Statistical score]
For a parametric family of densities $p_\theta(\vx)$ with $\theta\in\Theta\subset\Reals^d$,
the \term{statistical score} at $\vx$ is the gradient of the log-likelihood with
respect to the \emph{parameter},
\begin{equation}
  \vect s_\theta(\vx)=\nabla_\theta\log p_\theta(\vx)\in\Reals^d .
\end{equation}
\end{definition}

\begin{definition}[Diffusion score]
For a density $\Pt$ on $\Reals^{\nsites}$, the \term{diffusion score} (also Stein score)
at $\vx$ is the gradient of the log-density with respect to the \emph{data},
\begin{equation}
  \score(\vx,t)=\nabla_{\vx}\log\Pt(\vx)\in\Reals^{\nsites}.
\end{equation}
\end{definition}

Both are gradients of a log-density; they differ in the variable of
differentiation, in their dimension ($d$ against $\nsites$), and in their use. The
statistical score is a random vector whose covariance is the Fisher information and
whose sum over a sample is set to zero by maximum likelihood. The diffusion score is
a vector field that drives the reverse process and that a network is trained to
approximate. The paper writes ``the score'' for the diffusion score throughout, and
``the likelihood score'', ``per-sequence scores obtained by Fisher's identity'' or
``$\partial_\theta\log\Pt$'' for the statistical one. The two meet in one place: the
paper's Theorem 2 computes the Fisher information of a parameter of $\Pz$ from the
noised law $\Pt$, so the \emph{statistical} score $\partial_\theta\log\Pt(\vx)$ is
a function on the same space as the \emph{diffusion} score $\nabla_{\vx}\log\Pt(\vx)$.
A third, harmless usage, ``an arm is scored on the test bundle'', means evaluated,
and is avoided here.

\section{The score function and Fisher information}

Fix a family $p_\theta$, $\theta\in\Theta$ open in $\Reals^d$, and write
$\E_\theta$ for the expectation under $p_\theta$.

\begin{assumption}[Regularity]
\label{ass:regular}
(R1) The support $\{\vx:p_\theta(\vx)>0\}$ does not depend on $\theta$. (R2) $\theta\mapsto p_\theta(\vx)$
is continuously differentiable for every $\vx$. (R3) Differentiation with respect to
$\theta$ may be taken under the integral sign in $\int p_\theta\dd\vx$ and in
$\int T(\vx)p_\theta(\vx)\dd\vx$ for the statistics $T$ considered (a dominating
integrable bound on $\|\nabla_\theta p_\theta\|$ locally in $\theta$ suffices).
\end{assumption}

\begin{proposition}[The score has mean zero]
\label{prop:score-mean-zero}
Under \cref{ass:regular}, $\E_\theta[\vect s_\theta(\vx)]=\vect 0$.
\end{proposition}
\begin{proof}
$\E_\theta[\vect s_\theta]=\int\frac{\nabla_\theta p_\theta}{p_\theta}p_\theta\dd\vx=\int\nabla_\theta p_\theta\dd\vx=\nabla_\theta\int p_\theta\dd\vx=\nabla_\theta1=\vect0$,
using (R1) so that the integrand is defined on the support and (R3) to exchange
$\nabla_\theta$ and $\int$.
\end{proof}

\begin{definition}[Fisher information]
The \term{Fisher information matrix} of the family at $\theta$ is the covariance of
the score,
\begin{equation}
  \mat I(\theta)=\E_\theta\bigl[\vect s_\theta(\vx)\vect s_\theta(\vx)^{\!\top}\bigr],
  \qquad I_{jk}(\theta)=\E_\theta\bigl[\partial_j\log p_\theta\,\partial_k\log p_\theta\bigr].
  \label{eq:fisher}
\end{equation}
For a scalar parameter it is the number $I(\theta)=\E_\theta[(\partial_\theta\log p_\theta)^2]$.
\end{definition}

The paper's $I_t(\theta)=\E_{\Pt}[(\partial_\theta\log\Pt)^2]$ is this definition for
the family $\{\Pt^\theta\}$ of noised laws at fixed $t$, one sequence being one
observation.

\begin{proposition}[Four properties]
\label{prop:fisher-props}
Under \cref{ass:regular}, with second derivatives where stated:
\begin{enumerate}[label=(\roman*)]
\item (second-derivative form) if $\theta\mapsto p_\theta$ is twice differentiable and
  (R3) holds for the second derivatives, $I_{jk}(\theta)=-\E_\theta[\partial_j\partial_k\log p_\theta]$;
\item (additivity) for $\nseq$ independent observations with a common $\theta$ the
  information is $\nseq\,\mat I(\theta)$; for independent observations from different
  families it is the sum of their informations;
\item (reparametrisation) if $\theta=g(\psi)$ with $g$ differentiable and Jacobian
  $\mat J=\partial\theta/\partial\psi$, then $\mat I_\psi(\psi)=\mat J^{\!\top}\mat I_\theta(g(\psi))\mat J$;
  for scalars, $I_\psi=(g'(\psi))^2I_\theta$;
\item (data processing) if $\vect y=T(\vx)$ is a statistic, the information carried by
  $\vect y$ about $\theta$ is at most that carried by $\vx$, in the Loewner order.
\end{enumerate}
\end{proposition}
\begin{proof}
(i) Differentiate $\partial_j\log p=\partial_jp/p$ once more:
$\partial_j\partial_k\log p=\partial_j\partial_kp/p-(\partial_jp)(\partial_kp)/p^2$, and
$\E_\theta[\partial_j\partial_kp/p]=\int\partial_j\partial_kp\dd\vx=\partial_j\partial_k1=0$ by (R3);
the second term's expectation is $-I_{jk}$. (ii) The log-likelihood of independent
observations is a sum, so the score is a sum of independent mean-zero vectors and
its covariance is the sum of their covariances. (iii) By the chain rule
$\nabla_\psi\log p=\mat J^{\!\top}\nabla_\theta\log p$, and covariances transform as
$\mat J^{\!\top}\Cov\mat J$. (iv) Write $p_\theta(\vx)=p_\theta(\vect y)\,p_\theta(\vx\mid\vect y)$;
the score of $\vx$ is the score of $\vect y$ plus the score of the conditional law,
and the two are uncorrelated because the conditional score has conditional mean zero
given $\vect y$ (\cref{prop:score-mean-zero} applied to the conditional family); so
$\mat I_{\vx}=\mat I_{\vect y}+\E[\mat I_{\vx\mid\vect y}]\succeq\mat I_{\vect y}$.
\end{proof}

\why{Three of these are used by the paper without comment. Additivity is why the
bound for $\nseq$ sequences is $1/(\nseq I_t)$. Reparametrisation is how the paper's
Cram\'er--Rao reference for the excess kurtosis is obtained ``through the chain rule
at the Laplace member of the one-parameter shape family'': the information about
$\beta$ is computed, and the information about $\gamma_2=\kappa_4/\ivar^2$ as a function
of $\beta$ is $I_\beta/(\dd\gamma_2/\dd\beta)^2$, so the reference standard deviation of
an unbiased estimator of $\gamma_2$ is $|\dd\gamma_2/\dd\beta|/\sqrt{\nseq I_\beta}$.
Data processing is why the information of the noised sequence is at most that of the
clean one, and why the ``clean-data floor'' of the paper's Figure 3 is a floor.}

\section{The Cram\'er--Rao bound}

\begin{theorem}[Cram\'er--Rao, scalar]
\label{thm:cr-scalar}
Let $\theta$ be scalar, let \cref{ass:regular} hold, let $0<I(\theta)<\infty$, and let
$T(\vx)$ be an estimator with $\E_\theta[T^2]<\infty$ and bias $b(\theta)=\E_\theta[T]-\theta$
differentiable with (R3). Then
\begin{equation}
  \Var_\theta(T)\ \ge\ \frac{\bigl(1+b'(\theta)\bigr)^2}{I(\theta)},
  \qquad\text{and in particular}\qquad \Var_\theta(T)\ge\frac1{I(\theta)}\ \text{if $T$ is unbiased.}
  \label{eq:cr-scalar}
\end{equation}
\end{theorem}
\begin{proof}
Compute the covariance of $T$ with the score. Since the score has mean zero,
\[
  \Cov_\theta(T,s_\theta)=\E_\theta[T\,s_\theta]=\int T(\vx)\,\frac{\partial_\theta p_\theta(\vx)}{p_\theta(\vx)}\,p_\theta(\vx)\dd\vx
  =\partial_\theta\int T(\vx)p_\theta(\vx)\dd\vx=\partial_\theta\E_\theta[T]=1+b'(\theta),
\]
using (R3) for the exchange. The Cauchy--Schwarz inequality for covariances,
$\Cov(T,s)^2\le\Var(T)\Var(s)$, with $\Var_\theta(s_\theta)=I(\theta)$, gives
$(1+b'(\theta))^2\le\Var_\theta(T)\,I(\theta)$, which is \cref{eq:cr-scalar}.
\end{proof}

Equality in Cauchy--Schwarz holds if and only if $T-\E T$ is a multiple of the score,
$T(\vx)-\theta=c(\theta)\,s_\theta(\vx)$; an estimator that achieves the bound at every
$\theta$ is called \term{efficient}, and exists only for exponential families in
their mean parameterisation. The bound is a property of the family, not of any
estimator; that is why the paper calls it a \emph{reference} and not a target.

\begin{theorem}[Cram\'er--Rao, vector]
\label{thm:cr-vector}
Let $\theta\in\Reals^d$, let \cref{ass:regular} hold with $\mat I(\theta)$ finite and
non-singular, and let $\vect T(\vx)\in\Reals^d$ be unbiased for $\theta$ with finite
second moments. Then $\Cov_\theta(\vect T)\succeq\mat I(\theta)^{-1}$ in the Loewner order,
i.e. $\vect v^{\!\top}\Cov_\theta(\vect T)\vect v\ge\vect v^{\!\top}\mat I(\theta)^{-1}\vect v$
for every $\vect v$.
\end{theorem}
\begin{proof}
As in the scalar case, $\Cov_\theta(\vect T,\vect s_\theta)=\E_\theta[\vect T\vect s_\theta^{\!\top}]=\nabla_\theta\E_\theta[\vect T]^{\!\top}=\Id$.
The joint covariance of the stacked vector $(\vect T,\vect s_\theta)$ is
$\begin{psmallmatrix}\Cov(\vect T)&\Id\\ \Id&\mat I\end{psmallmatrix}\succeq0$, and a
block matrix $\begin{psmallmatrix}\mat A&\mat B\\ \mat B^{\!\top}&\mat C\end{psmallmatrix}$
with $\mat C\succ0$ is positive semidefinite if and only if its Schur complement
$\mat A-\mat B\mat C^{-1}\mat B^{\!\top}\succeq0$; here that is $\Cov(\vect T)-\mat I^{-1}\succeq0$.
\end{proof}

\section{Nuisance parameters and the efficient information}
\label{sec:efficient}

Partition $\theta=(\psi,\eta)$ into a coordinate of interest $\psi$ (scalar, say) and
nuisance parameters $\eta$, and partition the information accordingly,
$\mat I=\begin{psmallmatrix}I_{\psi\psi}&\mat I_{\psi\eta}\\ \mat I_{\eta\psi}&\mat I_{\eta\eta}\end{psmallmatrix}$.
By \cref{thm:cr-vector} the variance of an unbiased estimator of $\psi$ that must
also estimate $\eta$ is bounded below by $(\mat I^{-1})_{\psi\psi}$, and by the block
inverse formula
\begin{equation}
  (\mat I^{-1})_{\psi\psi}=\frac1{I^{\mathrm{eff}}_{\psi}},\qquad
  I^{\mathrm{eff}}_{\psi}=I_{\psi\psi}-\mat I_{\psi\eta}\mat I_{\eta\eta}^{-1}\mat I_{\eta\psi},
  \label{eq:efficient}
\end{equation}
the Schur complement. $I^{\mathrm{eff}}_\psi$ is the \term{efficient information}
about $\psi$: the information that survives after the nuisance has been estimated. It
is at most $I_{\psi\psi}$, with equality if and only if the scores of $\psi$ and of $\eta$
are uncorrelated. The ratio $I_{\psi\psi}/I^{\mathrm{eff}}_\psi\ge1$ is the
\term{nuisance cost}, the quantity \code{summarise} in
\code{src/markov\_score/information.py} reports. Equivalently, $I^{\mathrm{eff}}_\psi$
is the variance of the \term{efficient score}, the residual of the $\psi$-score after
its linear regression on the $\eta$-score.

\begin{proposition}[Nuisance at another Hermite degree]
\label{prop:nuisance-degree}
Suppose that, as $t\to\infty$ and with $\alpha=e^{-t}$, the scores of two coordinates
of the noised family have expansions
$s_\psi=\alpha^mA(\vx)+\alpha^{m+2}A_2(\vx)+\dots$ and $s_\eta=\alpha^{m'}B(\vx)+\alpha^{m'+2}B_2(\vx)+\dots$
with $A$, $B$ polynomials that are $\gamma$-orthogonal, $\E_\gamma[AB]=0$, and with
$I_{\psi\psi}=\alpha^{2m}\E_\gamma[A^2](1+O(\alpha^2))$, $I_{\eta\eta}=\alpha^{2m'}\E_\gamma[B^2](1+O(\alpha^2))$,
$\E_\gamma[B^2]>0$, and $\Pt=\gamma(1+O(\alpha^2))$ on the relevant moments. Then
$I_{\psi\eta}=O(\alpha^{m+m'+2})$ and
\begin{equation}
  I^{\mathrm{eff}}_\psi=I_{\psi\psi}\bigl(1+O(\alpha^4)\bigr)=\alpha^{2m}\E_\gamma[A^2]\bigl(1+O(\alpha^2)\bigr):
\end{equation}
the efficient information has the same rate and the same leading constant as the raw
information. If instead $m=m'$ and $\E_\gamma[AB]\ne0$, the correction is of the same
order as $I_{\psi\psi}$ and the efficient information is smaller by an $O(1)$ factor.
\end{proposition}
\begin{proof}
$I_{\psi\eta}=\E_{\Pt}[s_\psi s_\eta]=\alpha^{m+m'}\E_{\Pt}[AB]+O(\alpha^{m+m'+2})$, and
$\E_{\Pt}[AB]=\E_\gamma[AB]+O(\alpha^2)=O(\alpha^2)$. Hence
$I_{\psi\eta}^2/I_{\eta\eta}=O(\alpha^{2m+2m'+4})/\Theta(\alpha^{2m'})=O(\alpha^{2m+4})$, which is
$O(\alpha^4)$ relative to $I_{\psi\psi}=\Theta(\alpha^{2m})$. In the same-degree case the
cross term is $\alpha^{2m}\E_\gamma[AB]\ne0$ at leading order and the Schur complement
subtracts a fixed fraction.
\end{proof}

\why{This is the paper's sentence ``the leading orders of a degree-2 and a degree-4
coordinate do not couple, because $\Hen{2}$ and $\Hen{4}$ are $\gamma$-orthogonal, so the
efficient information of a shape coordinate after lower-degree nuisance coordinates
are estimated has the same rate and constant as its raw information; a nuisance of
the same degree costs an $O(1)$ factor''. For the generalised-Gaussian chain, $\psi=\beta$
has $m=4$ with $A$ a degree-four Hermite polynomial, and $\eta=(\corr,\ivar)$ have
$m'=2$ with $B$ a degree-two one. The proposition says the $e^{-8t}$ law and the constant
$\tfrac1{24}\Frob{\partial_\beta\kfour}^2$ of the paper's Theorem 2 hold for the efficient
information that the paper's Figure 2b actually plots, not only for the diagonal entry.
It also says why the eight-component mixture's shape coordinates, which are all of
degree four, cost each other an $O(1)$ factor, one reason the paper calls its
one-parameter reference a yardstick.}

The paper measures the information matrix of $(\corr,\ivar,\beta)$ by the
\term{outer-product estimator}: with per-sequence scores $\vect g_\mu=\nabla_\theta\log\Pt(\vx^\mu)$
at the true $\theta$, $\widehat{\mat I}=\tfrac1{\nseq}\sum_\mu\vect g_\mu\vect g_\mu^{\!\top}$,
which is unbiased for $\mat I$ by \cref{eq:fisher} (it is not the observed
information $-\nabla^2\log p$, which would need second derivatives). The function
\code{outer\_product\_information} of \code{src/markov\_score/information.py} forms it,
and \code{efficient\_information} inverts it coordinate by coordinate as in
\cref{eq:efficient}. Its own sampling error is what limits the paper's chain
verification of the shape constant; the scalar constant is verified by quadrature
instead (E21).

\section{The Gaussian family in closed form}
\label{sec:gaussian-fisher}

\begin{proposition}[Fisher information of a zero-mean Gaussian family]
\label{prop:gaussian-fisher}
Let $\vx\sim\normal{\vect0}{\Sig(\theta)}$ on $\Reals^{\nsites}$ with $\Sig(\theta)$
positive definite and differentiable in $\theta\in\Reals^d$. Then
\begin{equation}
  I_{jk}(\theta)=\tfrac12\tr\bigl(\Sig^{-1}\partial_j\Sig\,\Sig^{-1}\partial_k\Sig\bigr).
  \label{eq:gaussian-fisher}
\end{equation}
\end{proposition}
\begin{proof}
The log-density is $\log p_\theta(\vx)=-\tfrac{\nsites}2\log2\pi-\tfrac12\log\det\Sig-\tfrac12\vx^{\!\top}\Sig^{-1}\vx$.
Using $\partial_j\log\det\Sig=\tr(\Sig^{-1}\partial_j\Sig)$ and $\partial_j\Sig^{-1}=-\Sig^{-1}(\partial_j\Sig)\Sig^{-1}$,
\[
  \partial_j\log p_\theta(\vx)=-\tfrac12\tr(\Sig^{-1}\partial_j\Sig)+\tfrac12\vx^{\!\top}\Sig^{-1}(\partial_j\Sig)\Sig^{-1}\vx
  =\tfrac12\bigl(\vx^{\!\top}\mat B_j\vx-\E[\vx^{\!\top}\mat B_j\vx]\bigr),
\]
with $\mat B_j=\Sig^{-1}(\partial_j\Sig)\Sig^{-1}$ symmetric, since $\E[\vx^{\!\top}\mat B\vx]=\tr(\mat B\Sig)=\tr(\Sig^{-1}\partial_j\Sig)$.
The information is the covariance of these centred quadratic forms. For
$\vx\sim\normal{\vect0}{\Sig}$ and symmetric $\mat B$, $\mat C$, Isserlis' theorem
(chapter on Gaussian integrals and Hermite polynomials) gives
$\E[\vx^{\!\top}\mat B\vx\,\vx^{\!\top}\mat C\vx]=\tr(\mat B\Sig)\tr(\mat C\Sig)+2\tr(\mat B\Sig\mat C\Sig)$,
so $\Cov(\vx^{\!\top}\mat B\vx,\vx^{\!\top}\mat C\vx)=2\tr(\mat B\Sig\mat C\Sig)$. With the factors
$\tfrac12\cdot\tfrac12$, $I_{jk}=\tfrac14\cdot2\tr(\mat B_j\Sig\mat B_k\Sig)=\tfrac12\tr(\Sig^{-1}\partial_j\Sig\Sig^{-1}\partial_k\Sig)$.
\end{proof}

\paragraph{The noised Gaussian chain.}
For the Gaussian chain of the chapter on Markov chains, $\vx_t\sim\normal{\vect0}{\Sig_t(\theta)}$
exactly, with $\Sig_t=e^{-2t}\Sig_0(\theta)+\Dt\Id$ and $\theta=(\corr,\ivar)$, so
$\partial_\theta\Sig_t=e^{-2t}\partial_\theta\Sig_0$ and \cref{eq:gaussian-fisher} gives the
paper's closed form
\begin{equation}
  I_t(\theta)=\tfrac12\tr\bigl(\Sig_t^{-1}\partial_\theta\Sig_t\,\Sig_t^{-1}\partial_\theta\Sig_t\bigr)
  =\tfrac12e^{-4t}\tr\bigl(\Sig_t^{-1}\partial_\theta\Sig_0\,\Sig_t^{-1}\partial_\theta\Sig_0\bigr),
  \label{eq:gaussian-chain-fisher}
\end{equation}
exact at every $t$ and every $\corr$; this is the curve of the paper's Figure 2a and
the function \code{gaussian\_chain\_fisher} of \code{src/markov\_score/theory.py},
which evaluates $\partial_\theta\Sig_0$ by central differences. Its large-$t$ limit is
immediate: $\Sig_t\to\Id$, so
\begin{equation}
  e^{4t}I_t(\theta)\ \longrightarrow\ \tfrac12\tr\bigl(\partial_\theta\Sig_0\,\partial_\theta\Sig_0\bigr)=\tfrac12\Frob{\partial_\theta\Sig_0}^2,
\end{equation}
the constant of the paper's equation for a covariance coordinate, checked in Lean as
target 3 (\code{lean/README.md}). The approach to the limit is governed by the
spectrum of $\Sig_0$: writing $\Sig_t=\Id+e^{-2t}(\Sig_0-\Id)$, the Neumann series
$\Sig_t^{-1}=\Id-e^{-2t}(\Sig_0-\Id)+e^{-4t}(\Sig_0-\Id)^2-\dots$ converges in operator norm
exactly when $e^{-2t}\lmax(\Sig_0-\Id)<1$, i.e. $e^{2t}>\lmax(\Sig_0)-1$, and its first
correction is of relative size $e^{-2t}\lmax(\Sig_0)$: this is the origin of the paper's
$\epsilon_t=O(e^{-2t}\lmax(\Sig_0))$ and of the crossover $\tstar=\tfrac12\log\lmax(\Sig_0)$
below which the measured slope is shallower than $-4$. The same formula gives the
off-diagonal entry $I_{\corr\ivar}$, hence the correlation between the two scores and
the efficient information of each, which the paper's Figure 2a uses; the efficient
slope is already close to $-4$ below $\tstar$ where the raw diagonal is not, because
the nuisance coupling of $\corr$ with $\ivar$ is strong there.

\section{Local asymptotic normality and the H\'ajek--Le Cam theorem}
\label{sec:lan}

Fisher information bounds the variance of \emph{unbiased} estimators. Nothing in the
Cram\'er--Rao bound stops a biased estimator from having a smaller mean squared error
at some $\theta$ (the constant estimator has zero variance). The statement that
$1/I$ is nevertheless the right yardstick for \emph{every} estimator is asymptotic and
local, and it is the theorem the paper's corollary on sample complexity invokes.

\paragraph{In plain words.}
Take $\nseq$ independent sequences from $p_\theta$ and look at parameter values within
a $1/\sqrt{\nseq}$ neighbourhood of $\theta$, the scale on which the data can
distinguish them. As $\nseq\to\infty$ the whole experiment, restricted to that
neighbourhood, looks like observing one Gaussian vector with mean $\vect h$ (the
rescaled parameter offset) and covariance $\mat I(\theta)^{-1}$. In that Gaussian
experiment no estimator of $\vect h$ has worst-case risk below that of the observation
itself, whose covariance is $\mat I^{-1}$. So no estimator sequence can beat
$\mat I(\theta)^{-1}/\nseq$ uniformly over a shrinking neighbourhood, however biased,
and the maximum-likelihood estimator, which becomes that Gaussian observation,
attains it.

\begin{definition}[Local asymptotic normality]
The family is \term{locally asymptotically normal} (LAN) at $\theta$ with information
$\mat I(\theta)$ if, for every $\vect h\in\Reals^d$, the log-likelihood ratio of $\nseq$
independent observations satisfies
\begin{equation}
  \log\prod_{\mu=1}^{\nseq}\frac{p_{\theta+\vect h/\sqrt{\nseq}}(\vx^\mu)}{p_\theta(\vx^\mu)}
  =\vect h^{\!\top}\vect\Delta_{\nseq}-\tfrac12\vect h^{\!\top}\mat I(\theta)\vect h+o_{p}(1),
  \qquad \vect\Delta_{\nseq}=\frac1{\sqrt{\nseq}}\sum_\mu\vect s_\theta(\vx^\mu)\Rightarrow\normal{\vect0}{\mat I(\theta)},
\end{equation}
under $p_\theta$. A sufficient condition is that the family be
\term{differentiable in quadratic mean} at $\theta$: $\sqrt{p_{\theta+\vect h}}=\sqrt{p_\theta}+\tfrac12\vect h^{\!\top}\vect s_\theta\sqrt{p_\theta}+r_{\vect h}$
with $\int r_{\vect h}^2=o(|\vect h|^2)$; a continuously differentiable log-density with
continuous finite information suffices for that.
\end{definition}

\begin{theorem}[H\'ajek--Le Cam local asymptotic minimax]
\label{thm:lam}
Let the family be LAN at $\theta$ with non-singular $\mat I(\theta)$, let
$\ell:\Reals^d\to[0,\infty)$ be bowl-shaped (its sublevel sets convex and symmetric),
and let $\vect T_{\nseq}$ be any sequence of estimators. Then
\begin{equation}
  \lim_{c\to\infty}\ \liminf_{\nseq\to\infty}\ \sup_{|\vect h|\le c}\
  \E_{\theta+\vect h/\sqrt{\nseq}}\,\ell\bigl(\sqrt{\nseq}(\vect T_{\nseq}-\theta-\vect h/\sqrt{\nseq})\bigr)
  \ \ge\ \E\,\ell(\vect Z),\qquad \vect Z\sim\normal{\vect0}{\mat I(\theta)^{-1}} .
  \label{eq:lam}
\end{equation}
\end{theorem}

With the quadratic loss $\ell(\vect u)=(v^{\!\top}\vect u)^2$ for a coordinate direction,
$\E\ell(\vect Z)=(\mat I^{-1})_{\psi\psi}=1/I^{\mathrm{eff}}_\psi$, and \cref{eq:lam} reads:
for any estimator sequence $\hat\psi_{\nseq}$,
\begin{equation}
  \lim_{c\to\infty}\liminf_{\nseq\to\infty}\sup_{|\theta'-\theta|\le c/\sqrt{\nseq}}\nseq\,\E_{\theta'}(\hat\psi_{\nseq}-\psi')^2\ \ge\ \frac1{I^{\mathrm{eff}}_\psi(\theta)} .
  \label{eq:lam-quadratic}
\end{equation}
This is the statement the paper's corollary uses, with $I^{\mathrm{eff}}$ the efficient
information of the paper's general-order theorem at fixed $t$, so that the right-hand
side is $\Theta(e^{2mt})$ for a coordinate of lowest Hermite degree $m$. Two
features of the exact statement deserve attention. The supremum over a shrinking
neighbourhood is essential: without it a superefficient estimator can beat $1/I$ at
isolated points. And the outer limit $c\to\infty$ is part of the theorem; for a fixed
$c$ the local minimax risk of the limiting Gaussian experiment over $|\vect h|\le c$ is
strictly below $\E\ell(\vect Z)$ (a bounded normal-mean problem has a shrinkage
estimator with smaller worst-case risk), and the bound is approached as $c$ grows. The
paper's corollary writes the bound with a fixed $c$ and without the outer limit; read
literally that is slightly stronger than \cref{eq:lam-quadratic}, and the intended
statement is \cref{eq:lam-quadratic}.

\paragraph{What ``maximum likelihood attains it'' means.}
Under the regularity that gives LAN plus conditions for consistency, the
maximum-likelihood estimator satisfies $\sqrt{\nseq}(\hat\theta_{\nseq}-\theta)\Rightarrow\normal{\vect0}{\mat I(\theta)^{-1}}$
under $p_\theta$, locally uniformly, and is \term{regular} in the sense of the
convolution theorem; so its rescaled risk converges to $\E\ell(\vect Z)$ and equality
holds in \cref{eq:lam}. Three things this does \emph{not} say, all named by the paper.
(i) It is a statement at fixed $t$ as $\nseq\to\infty$; the LAN approximation needs
$\nseq\gg1/I_t$, i.e. $\nseq\gg e^{2mt}$, so the limits $\nseq\to\infty$ and $t\to\infty$
do not commute and nothing is claimed uniformly in $t$. (ii) It is a statement about
\emph{local} precision, on the scale $1/\sqrt{\nseq I_t}$; once that scale exceeds the
range of a bounded coordinate (the excess kurtosis of the mixture family is bounded
below by $-2$ and the paper's shape coordinate has range $3$), the neighbourhood
$|\theta'-\theta|\le c/\sqrt{\nseq}$ is no longer inside the parameter space and the
statement is empty; the paper reads its estimator's collapse at that level as the
information running out, not as a failure of the estimator. (iii) It says nothing
about an optimiser stopped at a finite budget, nor about an estimator whose iteration
count is selected on validation.

\section{Why ``at the bound'' is a yardstick and not an efficiency theorem}
\label{sec:yardstick}

The paper compares, level by level, the dispersion across sixteen replicates of its
fitted correlation and fitted excess kurtosis with the Cram\'er--Rao reference
$1/\sqrt{\nseq I^{\mathrm{eff}}_t}$. Four facts make that comparison a yardstick.

\begin{enumerate}
\item The reference is computed in a \emph{one-parameter} shape family (the
  generalised Gaussian at its Laplace member, with $(\corr,\ivar)$ as lower-degree
  nuisances that cost nothing to leading order, \cref{prop:nuisance-degree}). The
  estimator fits an \emph{eight-component mixture}, whose shape coordinates are all of
  degree four and cost each other $O(1)$ factors. The two families share the Laplace
  law as a point but not their tangent spaces, so the reference is not the bound for
  the estimator's own family.
\item The Cram\'er--Rao bound is for \emph{unbiased} estimators, and \cref{thm:cr-scalar}
  shows that a bias with $b'(\theta)<0$ lowers the bound by $(1+b')^2$. Early stopping
  selected on validation is a shrinkage device: it pulls the fitted shape toward the
  Gaussian innovation as the information falls. A dispersion \emph{below} the
  unbiased reference, which the paper reports at some levels, is therefore not a
  contradiction; it is the signature of shrinkage, and the paper reports bias and
  dispersion separately for that reason.
\item $\nseq=512$ is fixed and moderate, whereas the theorem is asymptotic. At the
  levels where $1/\sqrt{\nseq I_t}$ is a sizeable fraction of the coordinate's range,
  the local statement is empty (\cref{sec:lan}, point (ii)).
\item The reported dispersion is a sixteen-replicate standard deviation, whose own
  relative sampling error is about $1/\sqrt{2\cdot15}\approx18\%$ under normality
  (chapter on the statistics of the experiments); ``within the sampling error of a
  sixteen-replicate standard deviation'' is the most that can be claimed.
\end{enumerate}
The defensible statement, and the one the paper makes, is that at the low levels the
fits are unimodal and their dispersion is consistent with the one-parameter reference,
that collapse to a Gaussian innovation begins at a stated level, and that it is total
at the level where the reference crosses the coordinate's range.

\section{Fisher's identity for latent-variable models}
\label{sec:fisher-identity}

The paper's estimator maximises the marginal likelihood of the noised sequences,
$\Lik(\params)=\sum_\mu\log\Pt[t_\mu,\params](\vx^\mu)$, where each term is an integral
over the unobserved clean chain. Its gradient is obtained without differentiating
through the integral, by an identity due to Fisher.

\begin{proposition}[Fisher's identity]
\label{prop:fisher-identity}
Let $p_\theta(\va,\vx)$ be a joint density of a latent $\va$ and an observation $\vx$,
with marginal $p_\theta(\vx)=\int p_\theta(\va,\vx)\dd\va$, differentiable in $\theta$ with
(R3). Then
\begin{equation}
  \nabla_\theta\log p_\theta(\vx)=\E_\theta\bigl[\nabla_\theta\log p_\theta(\va,\vx)\,\big|\,\vx\bigr],
  \label{eq:fisher-identity}
\end{equation}
the posterior expectation of the complete-data score.
\end{proposition}
\begin{proof}
\[
  \nabla_\theta\log p_\theta(\vx)=\frac{\nabla_\theta\int p_\theta(\va,\vx)\dd\va}{p_\theta(\vx)}
  =\int\frac{\nabla_\theta p_\theta(\va,\vx)}{p_\theta(\va,\vx)}\cdot\frac{p_\theta(\va,\vx)}{p_\theta(\vx)}\dd\va
  =\int\nabla_\theta\log p_\theta(\va,\vx)\,p_\theta(\va\mid\vx)\dd\va ,
\]
using (R3) to move the gradient inside the integral and $p_\theta(\va\mid\vx)=p_\theta(\va,\vx)/p_\theta(\vx)$.
\end{proof}

Equivalently, with the EM auxiliary function $\Qfun(\theta\mid\theta')=\E_{\theta'}[\log p_\theta(\va,\vx)\mid\vx]$,
\cref{eq:fisher-identity} says $\nabla_\theta\Qfun(\theta\mid\theta')|_{\theta=\theta'}=\nabla_\theta\log p_{\theta'}(\vx)$:
the gradient of the EM surrogate at the current point is the gradient of the marginal
log-likelihood, which is why EM's ascent directions are likelihood ascent directions
(chapter on maximum likelihood and EM).

\paragraph{For the chain.}
The complete-data density is \cref{eq:markov-joint} of the chapter on Markov chains
times the channel, $p_\theta(\va,\vx)=\mu(a_1)\prod_{i\ge2}\kernelth(a_i\mid a_{i-1})\prod_i\like(x_i\mid a_i)$,
and only the transition factors depend on $\theta$. So
\begin{equation}
\begin{aligned}
  \nabla_\theta\log p_\theta(\vx)&=\sum_{i\ge2}\E_\theta\bigl[\nabla_\theta\log\kernelth(a_i\mid a_{i-1})\,\big|\,\vx\bigr]\\
  &=\sum_{i\ge2}\iint\nabla_\theta\log\kernelth(a'\mid a)\,p_\theta(a_i=a',a_{i-1}=a\mid\vx)\dd a\dd a' ,
\end{aligned}
\end{equation}
a sum of integrals against the \term{pairwise posterior marginals} of consecutive
sites, which the forward--backward recursion returns in the same pass that computes
the site marginals. On the grid, with $\Xis_{kl}=\sum_\mu\sum_{i\ge2}p(a_i=\gridpt{k},a_{i-1}=\gridpt{l}\mid\vx^\mu)$
the expected transition mass (child first, parent second), the gradient of the whole
marginal log-likelihood is the inner product $\langle\Xis,\nabla_\theta\log\kernelth\rangle$
of the paper's equation (4). The same $\Xis$ is all the M-step of EM needs. The
identity is verified in the paper against finite differences of the grid evidence to
$10^{-9}$, and, independently of the recursion, \code{tests/test\_fisher\_bruteforce.py}
enumerates every latent path on a five-point grid and checks the evidence, the
orientation of $\Xis$ (a transposed $\Xis$ paired with a transposed kernel would pass a
finite-difference check), and the gradient against a finite difference of the
enumerated evidence. The per-sequence scores $\vect g_\mu$ that feed the outer-product
information estimator of \cref{sec:efficient} are computed by the batched version of
the same identity in \code{src/markov\_score/information.py}
(\code{per\_chain\_scores}), one term per edge, without materialising a per-sequence
$\ngrid\times\ngrid$ matrix.

\section{Worked example: the variance of one noised Gaussian coordinate}

Let $x\sim\normal{0}{\theta}$, one scalar observation, $\theta>0$. The log-density is
$-\tfrac12\log(2\pi\theta)-x^2/(2\theta)$, the score is $-1/(2\theta)+x^2/(2\theta^2)$, and
since $\E x^2=\theta$ and $\Var(x^2)=2\theta^2$,
\[
  I(\theta)=\Var\Bigl(\frac{x^2}{2\theta^2}\Bigr)=\frac{2\theta^2}{4\theta^4}=\frac1{2\theta^2},
\]
which agrees with \cref{eq:gaussian-fisher}: $\tfrac12\tr(\theta^{-1}\cdot1\cdot\theta^{-1}\cdot1)=1/(2\theta^2)$.
The Cram\'er--Rao bound for $\nseq$ unbiased observations is $2\theta^2/\nseq$, attained
by the sample second moment $\tfrac1{\nseq}\sum x_\mu^2$, whose variance is exactly
$2\theta^2/\nseq$: an efficient estimator, as the exponential-family remark after
\cref{thm:cr-scalar} predicts.

Now noise the coordinate: $x_t=e^{-t}a+\sqrt{\Dt}z$ with $a\sim\normal0\theta$, so
$x_t\sim\normal{0}{v_t}$ with $v_t=e^{-2t}\theta+1-e^{-2t}$ and $\partial_\theta v_t=e^{-2t}$.
By \cref{eq:gaussian-fisher},
\[
  I_t(\theta)=\frac{(\partial_\theta v_t)^2}{2v_t^2}=\frac{e^{-4t}}{2\bigl(1+e^{-2t}(\theta-1)\bigr)^2}.
\]
At $\theta=1$ the denominator is exactly $1$ at every $t$, so $I_t(1)=\tfrac12e^{-4t}$
\emph{exactly}: the constant $\tfrac12\Frob{\partial_\theta\Sig_0}^2=\tfrac12$ of the
paper's covariance-coordinate law, with no correction, because a unit-variance
coordinate keeps unit variance under the channel. This is the ``$I_te^{4t}$ equal to
one half'' that the paper's Appendix E checks for the variance coordinate of the scalar
law, and that \code{tests/test\_theory\_hermite.py} checks by quadrature for the
Laplace law as well (the paper's Theorem 2 says the constant is the same for every
law with an exponential moment). At $\theta\ne1$ the factor $(1+e^{-2t}(\theta-1))^{-2}$ is
the scalar version of the crossover: for $\theta=5$ the local log-slope of $I_t$ is
$-4+4\cdot4e^{-2t}/(1+4e^{-2t})$, about $-1.6$ at $t=0.5$, $-3.3$ at $t=1.5$ and $-4.0$ at
$t=3$, with the crossover $\tfrac12\log(\theta-1)=\tfrac12\log4\approx0.69$ at which the
correction term equals one. With $\nseq=512$ sequences at $t=1$ and $\theta=1$, the
reference standard deviation is $\sqrt{2e^{4}/512}\approx0.46$, more than the range of
a coordinate confined to an interval of length one: at that level the local theory
has nothing to say about such a coordinate, which is the content of point (ii) in
\cref{sec:lan}.

\where{{\sloppy The paper's Section 3 (the definition of $I_t(\theta)$, the closed-form Gaussian
information, the efficient-information sentence, Theorem 3 on the general order,
Corollary 4 and the three caveats after it), Section 4 (the Cram\'er--Rao reference,
``a yardstick, not an efficiency theorem'', Fisher's identity and equation (4)),
Appendix A (the proof of the corollary and the nuisance-projection argument),
Appendix C (the finite-difference verification) and Appendix E (E21, the scalar
constants). Experiments E01, E06, E12b, E19, E21. Code:
\code{src/markov\_score/theory.py} (\code{gaussian\_chain\_fisher},
\code{correlation\_information\_constant}, \code{shape\_information\_constant}),
\code{src/markov\_score/information.py} (\code{per\_chain\_scores},
\code{outer\_product\_information}, \code{efficient\_information}, \code{summarise}),
\code{src/markov\_score/shape\_kernel.py} (the shape coordinate and its gradient),
\code{tests/test\_fisher\_bruteforce.py}, \code{tests/test\_theory\_hermite.py},
\code{tests/test\_shape\_kernel.py}.\par}}

\paragraph{References.}
H. Cram\'er (1946), \emph{Mathematical Methods of Statistics}, Princeton University
Press (Chapter 32: the information inequality).
C. R. Rao (1945), ``Information and the accuracy attainable in the estimation of
statistical parameters'', Bulletin of the Calcutta Mathematical Society 37, 81--91.
E. L. Lehmann and G. Casella (1998), \emph{Theory of Point Estimation}, 2nd edition,
Springer (Chapter 2: Fisher information, nuisance parameters, efficiency).
A. W. van der Vaart (1998), \emph{Asymptotic Statistics}, Cambridge University Press
(Chapters 7 and 8: differentiability in quadratic mean, local asymptotic normality, the
convolution and local asymptotic minimax theorems).
J. H\'ajek (1972), ``Local asymptotic minimax and admissibility in estimation'',
Proceedings of the Sixth Berkeley Symposium on Mathematical Statistics and
Probability 1, 175--194.
O. Capp\'e, E. Moulines and T. Ryd\'en (2005), \emph{Inference in Hidden Markov Models},
Springer (Chapter 10: Fisher's identity, Louis' identity and the EM algorithm for
state-space models).
